\documentclass[11pt]{amsart}
\usepackage{amsmath,amssymb,amsthm,mathrsfs}
\usepackage{graphicx}
\usepackage[hidelinks]{hyperref}

\newtheorem{theorem}{Theorem}
\newtheorem{lemma}[theorem]{Lemma}

\newcommand{\N}{\mathbb N}
\newcommand{\R}{\mathbb R}
\newcommand{\U}{\mathcal U}
\newcommand{\cU}{\mathfrak c}
\newcommand{\Linfty}{L^\infty}
\newcommand{\ellinf}{\ell_\infty}

\title{A Lifted Dual Fiber Need Not Fill the Fiber Dual}
\author{Agent lane 12}
\date{August 9, 2026}

\begin{document}
\maketitle

\section*{Source and target}

Cakovi\'c--Lu\v{c}i\'c--Pasqualetto construct, for a complete probability
space $(X,\Sigma,\mu)$, a von Neumann lifting $\ell$, an
$L^\infty(\mu)$-Banach module $\mathscr M$, and $x\in X$, a canonical
isometric embedding
\[
  {\sf j}_x:\ell(\mathscr M^*)_x\longrightarrow(\ell\mathscr M_x)'.
\]
Immediately after Proposition 3.22 of \cite{CLP}, they ask whether
${\sf j}_x$ must be surjective, referring back to Problem 3.11 of
Di Marino--Lu\v{c}i\'c--Pasqualetto \cite{DMLP}.

\begin{center}
  \includegraphics[width=0.98\linewidth]{figures/open_problem_crop.png}
\end{center}

\section*{Result}

The answer is no.  A single null point, equipped by the lifting with a free
ultrafilter value, turns the lifted fiber of a constant $c_0$-bundle into the
Banach ultrapower $(c_0)_\U$.  The corresponding lifted dual fiber is only
$(\ell_1)_\U$.  The canonical map is the usual pairing map
\[
  (\ell_1)_\U\hookrightarrow ((c_0)_\U)',
\]
and a cardinality argument shows that it cannot be onto.

\section*{Proof intuition}

Put positive mass on every integer and zero mass on one extra point
$\infty$.  Modulo null sets, a bounded scalar function is just a bounded
sequence.  At $\infty$, choose its representative by taking a free
ultrafilter limit.  A bounded sequence of vectors in $c_0$ is then represented
at $\infty$ by its ultrapower class, so the lifted module fiber is
$(c_0)_\U$.  Module functionals act independently on the positive atoms; they
are precisely bounded sequences of elements of $c_0'=\ell_1$, and their
lifted fiber is $(\ell_1)_\U$.

The missing functionals are forced by size.  The space $(\ell_1)_\U$ has at
most continuum many elements.  In contrast, $(c_0)_\U$ contains an isometric
copy of $\ell_\infty$, obtained by taking longer and longer initial
truncations.  The $2^{\mathfrak c}$ ultrafilter-limit functionals on
$\ell_\infty$ extend by Hahn--Banach to $(c_0)_\U$.  Thus the latter dual has
strictly more elements than the image of the canonical map.

\begin{theorem}\label{thm:counterexample}
There exist a complete probability space $(X,\Sigma,\mu)$, a von Neumann
lifting $\ell$, a countably generated $L^\infty(\mu)$-Banach module
$\mathscr M$ represented by the constant separable fiber $c_0$, and a point
$x\in X$ such that
\[
  {\sf j}_x:\ell(\mathscr M^*)_x\longrightarrow(\ell\mathscr M_x)'
\]
is not surjective.
\end{theorem}

\begin{proof}
We work over the real scalars; the same construction works over the complex
scalars.  Let
\[
 X=\N\cup\{\infty\},\qquad \Sigma=\mathcal P(X),\qquad
 \mu(\{n\})=2^{-n}\ (n\geq1),\qquad \mu(\{\infty\})=0.
\]
This is a complete probability space.  Fix a free ultrafilter $\U$ on $\N$.
Every member of $L^\infty(\mu)$ has a unique bounded sequence of values on
$\N$, so define
\begin{equation}\label{eq:scalar-lifting}
 \ell a(n)=a_n\quad(n\in\N),
 \qquad
 \ell a(\infty)=\lim_{\U}a_n.
\end{equation}
Ultrafilter limits preserve sums, products, constants, order, and absolute
values.  Hence \eqref{eq:scalar-lifting} is a positive unital multiplicative
linear right inverse to the almost-everywhere quotient: it is a von Neumann
lifting.

Take
\[
  \mathscr M=L^\infty(\mu;c_0)=\ell_\infty(\N;c_0).
\]
We now identify its lifted module without invoking any unstated fiber formula.
Let
\[
 E_n=c_0\quad(n\in\N),
 \qquad E_\infty=(c_0)_\U,
\]
Let $\mathscr E$ be the module, over the bounded measurable scalar functions,
of all bounded sections $w=(w(x))_{x\in X}$ of this family, with pointwise
operations and pointwise norm.  It is an
$\mathcal L^\infty(\Sigma)$-Banach module.  Define
\begin{equation}\label{eq:vector-lifting}
 Lv(n)=v_n,
 \qquad Lv(\infty)=[v_n]_\U
 \quad\text{for }v=(v_n)\in\ell_\infty(\N;c_0).
\end{equation}
Then $L$ is linear and
\[
  |Lv|(n)=\|v_n\|,
  \qquad |Lv|(\infty)=\lim_\U\|v_n\|,
\]
so $|Lv|=\ell|v|$.  It also respects scalar multiplication: if
$a=(a_n)\in L^\infty(\mu)$, then
$[a_nv_n]_\U=(\lim_\U a_n)[v_n]_\U$, because $(v_n)$ is bounded.

The image of $L$ generates all of $\mathscr E$, in fact without taking a
closure.  Indeed, given $w\in\mathscr E$, put $v_n=w(n)$.  Choose a bounded
representative $(z_n)$ of
\[
  w(\infty)-[v_n]_\U\in(c_0)_\U.
\]
Then, as sections on $X$,
\[
  w=Lv+\mathbf 1_{\{\infty\}}Lz.
\]
Thus $(\mathscr E,L)$ satisfies exactly the defining characterization of the
module lifting, and uniqueness of liftings yields
\begin{equation}\label{eq:first-fiber}
  \ell\mathscr M_\infty\cong(c_0)_\U.
\end{equation}
For completeness, the lifted atoms here are the singletons $\{n\}$,
$n\in\N$: a free ultrafilter contains no singleton.  Hence $\infty$ lies
outside their union, and the fiber seminorm at $\infty$ is simply evaluation
of the pointwise norm there.  This verifies directly that the fiber of the
model above is $E_\infty$.

Next identify the module dual.  If $T\in\mathscr M^*$, then
$L^\infty(\mu)$-linearity implies that the $n$th coordinate of $T(v)$ depends
only on $v_n$.  Consequently there are $\phi_n\in c_0'=\ell_1$ such that
\[
  T(v)_n=\phi_n(v_n).
\]
The pointwise boundedness condition for a module functional gives
\[
  \sup_{n\in\N}\|\phi_n\|<\infty.
\]
Conversely, every bounded sequence
$(\phi_n)$ in $\ell_1$ defines such a module functional.  Therefore,
isometrically as modules,
\[
  \mathscr M^*=\ell_\infty(\N;\ell_1).
\]
Repeating the explicit lifting construction above with $c_0$ replaced by
$\ell_1$ gives
\begin{equation}\label{eq:dual-fiber}
  \ell(\mathscr M^*)_\infty\cong(\ell_1)_\U.
\end{equation}

Under \eqref{eq:first-fiber} and \eqref{eq:dual-fiber}, the defining identity
for ${\sf j}_\infty$ says that
\begin{equation}\label{eq:canonical-map}
 {\sf j}_\infty([\phi_n]_\U)([v_n]_\U)
   =\lim_\U\phi_n(v_n).
\end{equation}
This is the standard canonical isometric embedding
$(\ell_1)_\U\hookrightarrow((c_0)_\U)'$.

It remains to prove that \eqref{eq:canonical-map} is not onto.  Write
$\mathfrak c=2^{\aleph_0}$.  Since $|\ell_1|=\mathfrak c$,
\[
  |(\ell_1)_\U|leq |(\ell_1)^\N|
  =\mathfrak c^{\aleph_0}=\mathfrak c.
\]
On the other hand, $(c_0)_\U$ contains an isometric copy of $\ell_\infty$.
For $a=(a_k)_{k\geq1}\in\ell_\infty$, let
\[
  x_n(a)=(a_1,\ldots,a_n,0,0,\ldots)\in c_0,
  \qquad J(a)=[x_n(a)]_\U.
\]
Because $\U$ is free, every cofinite set belongs to $\U$, and hence
\[
 \|J(a)\|
 =\lim_\U\max_{1\leq k\leq n}|a_k|
 =\sup_{k\geq1}|a_k|.
\]
Thus $J:\ell_\infty\to(c_0)_\U$ is a linear isometry.

There are $2^{\mathfrak c}$ ultrafilters on $\N$.  Distinct ultrafilters
$\mathcal V$ give distinct norm-one functionals
\[
  q_{\mathcal V}(a)=\lim_{\mathcal V}a_n
  \qquad(a\in\ell_\infty),
\]
as is seen on characteristic sequences.  By Hahn--Banach, each
$q_{\mathcal V}\circ J^{-1}$ on $J(\ell_\infty)$ extends to a functional on
$(c_0)_\U$; extensions coming from distinct restrictions remain distinct.
Therefore
\[
  |((c_0)_\U)'|\geq2^{\mathfrak c}
  >\mathfrak c\geq|{\sf j}_\infty((\ell_1)_\U)|.
\]
The canonical isometric embedding ${\sf j}_\infty$ is not surjective.
\end{proof}

\section*{Compatibility and scope}

The source paper proves that the image of ${\sf j}_x$ is weak-$*$ dense.
There is no conflict: the proper subspace
${\sf j}_\infty((\ell_1)_\U)$ separates the points of $(c_0)_\U$, so it is
indeed weak-$*$ dense in the full dual.  The counterexample only refutes
unrestricted surjectivity.  It leaves intact possible positive theorems for
Hilbertian, reflexive, finite-dimensional, or otherwise restricted modules.

The example is already a constant measurable Banach bundle with separable
fiber $c_0$, and its section module is countably generated.  The pathology is
concentrated at the null point $\infty$; at every positive atom $n$, the
lifted fiber is just $c_0$ and the corresponding map is the familiar
$\ell_1\cong c_0'$.

\section*{Verification notes}

The proof is self-contained apart from three standard choice-based facts:
the existence of a free ultrafilter on $\N$, the fact that $\N$ has
$2^{\mathfrak c}$ ultrafilters, and Hahn--Banach.  The same strength of choice
is already inherent in the general lifting framework.  The two identifications
of lifted fibers were checked directly from the source's defining
characterization: preservation of the pointwise norm and generation by the
lifted image.

The highest-value review points are: (i) the generation identity
$w=Lv+\mathbf1_{\{\infty\}}Lz$; (ii) the coordinatewise description of the
module dual; and (iii) the standard count $2^{\mathfrak c}$ of ultrafilters
on a countable set.  Together they make the cardinal obstruction decisive.

\section*{Novelty and search bounds}

Before promotion, the run's registry, solution, attempt, and proof-gap
indexes were searched for arXiv:2508.19635 and for the phrases ``lifted dual
fiber'', ``$j_x$ surjective'', ``module lifting'', and ``Banach ultrapower''.
No duplicate result was found.  Bounded external searches on August 9, 2026,
used the exact wording of Problem 3.11 and combinations of
``lifted module dual fiber'', ``surjective'', ``ultrapower'', and
``$(c_0)_\mathcal U$''.  They recovered the original problem and the 2025
paper that repeats its open status, but no paper answering it.  Novelty should
still be treated as medium confidence pending expert bibliographic review.

\begin{thebibliography}{9}
\bibitem{CLP}
M. Cakovi\'c, D. Lu\v{c}i\'c, and E. Pasqualetto,
\emph{Tensor products of measurable Banach bundles},
arXiv:2508.19635, 2025.

\bibitem{DMLP}
S. Di Marino, D. Lu\v{c}i\'c, and E. Pasqualetto,
\emph{Representation theorems for normed modules},
Rev. Real Acad. Cienc. Exactas Fis. Nat. Ser. A-Mat. 119 (2025), article 19.
\end{thebibliography}

\end{document}
